% Источники: raw/IMG_7616.jpeg, raw/IMG_7617.jpeg (в этом порядке).
% Компиляция: tectonic 05.tex
% Продолжение: raw/IMG_7618.jpeg -- raw/IMG_7624.jpeg.
% Сокращения раскрыты; существенные исправления отмечены явно.
\ifdefined\lessonpreambleloaded
  \expandafter\endinput
\fi
\let\lessonpreambleloaded\empty
\documentclass[a4paper,11pt]{article}
\usepackage{fontspec}
\setmainfont{Times New Roman}
\usepackage[russian]{babel}
\usepackage{amsmath,amssymb}
\usepackage[margin=22mm]{geometry}
\usepackage{indentfirst}
\usepackage{needspace}
\usepackage{tikz}
\usetikzlibrary{patterns}
\usepackage[hidelinks]{hyperref}
\usepackage{fancyhdr}
\pagestyle{fancy}
\fancyhf{}
\renewcommand{\headrulewidth}{0pt}
\fancyfoot[C]{\footnotesize\url{https://github.com/bmstu-software-engineering/cs-s5-tv-26y}}
\fancyfoot[R]{\thepage}
\setlength{\emergencystretch}{2em}
\newcommand{\sourcepage}[1]{\par\Needspace{6\baselineskip}\medskip\noindent{\footnotesize\itshape Источник: \detokenize{#1}.}\par\nopagebreak\smallskip}
\newcommand{\editorial}[1]{\par\smallskip\noindent{\small\itshape Примечание к переносу. #1}\par\smallskip}
\newcommand{\lessontitle}[3]{\begin{center}{\large\bfseries Теория вероятностей. #1}\\[4pt]{\large #2}\\[4pt]#3\end{center}}
\newcommand{\parttitle}[1]{\par\Needspace{4\baselineskip}\medskip\noindent\textbf{#1}\par\nopagebreak\smallskip}
\tikzset{element/.style={draw,rectangle,minimum width=7mm,minimum height=5mm,fill=white},every picture/.style={line width=.5pt}}


\begin{document}

\begin{center}
  {\large\bfseries Теория вероятностей. Семинар №\,5}\\[4pt]
  {\large Классическое определение вероятности}\\[4pt]
  03.10.26
\end{center}

\section*{Обозначения}

$\Omega$~--- пространство равновозможных элементарных исходов.

$N$~--- мощность пространства $\Omega$, то есть число его элементов:
\[
  N=|\Omega|.
\]

$A$~--- случайное событие; событие является подмножеством
пространства исходов: $A\subseteq\Omega$.

$N_A$~--- мощность множества $A$:
\[
  N_A=|A|.
\]

\subsection*{Оформление}

$(x_1,x_2,\ldots,x_n)$~--- упорядоченный набор (кортеж): порядок важен.

$\{x_1,x_2,\ldots,x_n\}$~--- неупорядоченный набор (множество).

\section*{Задача 1. Два броска игральной кости}

Игральную кость бросают два раза. Найти вероятность того, что:
\begin{enumerate}
  \item выпадут одинаковые числа;
  \item сумма выпавших очков будет равна $8$.
\end{enumerate}

\textbf{Исход:} $(x_1,x_2)$; порядок важен.
Здесь $x_i\in\{1,2,3,4,5,6\}$~--- число очков,
выпавших при $i$-м броске, $i=1,2$.

\[
  \Omega=\left\{
  \begin{array}{cccc}
    (1,1), & (1,2), & \ldots, & (1,6),\\
    (2,1), & (2,2), & \ldots, & (2,6),\\
    \vdots & \vdots & \ddots & \vdots\\
    (6,1), & (6,2), & \ldots, & (6,6)
  \end{array}
  \right\}.
\]
\[
  N=|\Omega|=36.
\]

\newpage
% Начало raw/IMG_7617.jpeg: продолжение задачи 1.

\subsection*{1. Одинаковые числа}

Пусть $A$~--- событие «выпали одинаковые числа».
\[
  A=\{(1,1),(2,2),\ldots,(6,6)\},
  \qquad N_A=|A|=6.
\]
\[
  P(A)=\frac{N_A}{N}=\frac{6}{36}=\frac{1}{6}.
\]

\subsection*{2. Сумма очков равна восьми}

Пусть $B$~--- событие «сумма выпавших очков равна $8$».
\[
  B=\{(2,6),(3,5),(4,4),(5,3),(6,2)\},
  \qquad N_B=|B|=5.
\]
\[
  P(B)=\frac{N_B}{N}=\frac{5}{36}.
\]

\textbf{Ответ:} $P(A)=\dfrac{1}{6}$, $P(B)=\dfrac{5}{36}$.

\section*{Задача 2. Кость домино}

Наудачу выбирают кость домино. Найти вероятность того, что:
\begin{enumerate}
  \item выбранная кость~--- дубль;
  \item сумма очков равна $7$.
\end{enumerate}

\textbf{Исход:} $\{m,n\}$; порядок не важен.
Здесь $m,n\in\{0,1,\ldots,6\}$~--- числа очков на половинах кости.
Для записи выбираем $m\leq n$.
% Фигурные скобки сохранены как в исходнике: это обозначение кости,
% в том числе дубля {m,m}, а не буквальное множество чисел.
\[
  \Omega=\left\{
  \begin{array}{ccccc}
    \{0,0\}, & \{0,1\}, & \{0,2\}, & \ldots, & \{0,6\},\\
             & \{1,1\}, & \{1,2\}, & \ldots, & \{1,6\},\\
             &           & \ddots   &         & \vdots\\
             &           &          &         & \{6,6\}
  \end{array}
  \right\}.
\]

\sourcepage{IMG_7618.jpeg}
\[
 N=|\Omega|=28.
\]
\parttitle{а) Дубль}
Пусть $A$ --- событие «выбран дубль».
\[
 A=\bigl\{\{0,0\},\{1,1\},\{2,2\},\{3,3\},\{4,4\},\{5,5\},\{6,6\}\bigr\},
 \qquad N_A=|A|=7.
\]
\[
 P(A)=\frac{N_A}{N}=\frac7{28}=\frac14.
\]
\parttitle{б) Сумма очков равна семи}
\[
 B=\bigl\{\{1,6\},\{2,5\},\{3,4\}\bigr\},\qquad N_B=|B|=3,
 \qquad P(B)=\frac{N_B}{N}=\frac3{28}.
\]
\textbf{Ответ:} $P(A)=\dfrac14$, $P(B)=\dfrac3{28}$.

\section*{Основная формула комбинаторики}
Пусть имеется $m$ классов, в $i$-м классе $n_i$ элементов,
$i=1,\ldots,m$. Число способов выбрать по одному элементу каждого класса:
\[
 \boxed{N=n_1n_2\cdots n_m.}
\]
\textbf{Пример.} Код имеет вид «цифра--буква--цифра».
Для цифры имеется $10$ вариантов, для буквы --- $33$. Поэтому
\[
 N=10\cdot33\cdot10=3300.
\]
При подсчёте исходов различают сочетания и размещения, а также
выбор с повторениями и без повторений.

\sourcepage{IMG_7619.jpeg}
\subsection*{Сочетания и размещения}
\begin{center}
\renewcommand{\arraystretch}{1.8}
\begin{tabular}{|l|c|c|}
\hline
 & С повторениями & Без повторений \\
\hline
Порядок важен & $\widetilde A_n^m=n^m$ & $A_n^m=\dfrac{n!}{(n-m)!}$ \\
\hline
Порядок не важен & $\widetilde C_n^m=C_{n+m-1}^m$ & $C_n^m=\dfrac{n!}{m!(n-m)!}$ \\
\hline
\end{tabular}
\end{center}
$A$ обозначает размещения, $C$ --- сочетания.
Пример сочетаний без повторений: выбор $m$ конфет из $n$ разных.
Сочетания с повторениями в исходнике сопровождаются примером с картами.
\editorial{В рукописной таблице две формулы сочетаний стоят под противоположными заголовками.
Здесь их расположение исправлено.}

\section*{Задача 3. Шары в урне}
В урне $10$ шаров: $4$ чёрных и $6$ белых. Случайным образом
вынимают $3$ шара. Найти вероятности того, что:
\begin{enumerate}
 \item среди вынутых шаров нет чёрных;
 \item среди вынутых шаров ровно один чёрный.
\end{enumerate}
Пронумеруем чёрные шары числами $1,2,3,4$, белые --- $5,6,\ldots,10$.
Исход --- неупорядоченная тройка различных номеров:
\[
 \Omega=\bigl\{\{x_1,x_2,x_3\}:x_i\in\{1,\ldots,10\},\;
 x_i\ne x_j\text{ при }i\ne j\bigr\}.
\]
Это сочетания без повторений из $10$ по $3$.

\sourcepage{IMG_7620.jpeg}
\[
 N=|\Omega|=C_{10}^3=\frac{10!}{3!7!}.
\]
\parttitle{а) Нет чёрных шаров}
\[
 A=\bigl\{\{x_1,x_2,x_3\}:x_i\in\{5,\ldots,10\},\;
 x_i\ne x_j\text{ при }i\ne j\bigr\}.
\]
\[
 N_A=C_6^3=\frac{6!}{3!3!},\qquad
 P(A)=\frac{C_6^3}{C_{10}^3}=\frac16.
\]
\parttitle{б) Ровно один чёрный шар}
Один номер выбираем из $\{1,2,3,4\}$, два различных номера ---
из $\{5,\ldots,10\}$. По основной формуле комбинаторики
\[
 N_B=C_4^1C_6^2=\frac{4!}{1!3!}\frac{6!}{2!4!}
 =\frac{6!}{2!3!},
\]
\[
 P(B)=\frac{C_4^1C_6^2}{C_{10}^3}
 =\frac{3\cdot4\cdot5\cdot6}{8\cdot9\cdot10}=\frac12.
\]
\parttitle{Дополнительный пункт: три чёрных шара}
Внизу страницы записано событие $C$ --- «вынуты три чёрных шара»:
\[
 P(C)=\frac{C_4^3}{C_{10}^3}.
\]
\editorial{Дальнейшее вычисление этого пункта в исходнике не закончено.}

\sourcepage{IMG_7621.jpeg}
\section*{Задача 4. Лифт}
В семиэтажном доме на первом этаже в лифт вошли три человека.
Каждый может выйти на любом этаже, начиная со второго. Рассматриваются события:
\begin{enumerate}
 \item все выйдут на четвёртом этаже;
 \item все выйдут на одном этаже;
 \item все выйдут на разных этажах.
\end{enumerate}
Исход $(x_1,x_2,x_3)$, где $x_i\in\{2,3,\ldots,7\}$ --- этаж,
на котором выходит $i$-й человек. Порядок важен, повторения допустимы:
\[
 \Omega=\{(x_1,x_2,x_3):x_i\in\{2,\ldots,7\},\ i=1,2,3\},
 \qquad N=\widetilde A_6^3=6^3.
\]
Для третьего события $C$ этажи попарно различны, поэтому
\[
 N_C=A_6^3=\frac{6!}{3!},\qquad
 P(C)=\frac{6!}{3!6^3}=\frac{4\cdot5}{36}.
\]
\editorial{Решения первых двух пунктов на фотографии отсутствуют.}

\sourcepage{IMG_7622.jpeg}
\parttitle{Домашнее задание}
Вверху страницы указаны номера: 2.17, 2.23, 2.29, 2.35.

\section*{Задача 2.17. Буквы}
В исходнике приведены буквы слова «ЛИТЕРА» с номерами:
\[
 \begin{array}{c|cccccc}
 \text{буква}&\text{Л}&\text{И}&\text{Т}&\text{Е}&\text{Р}&\text{А}\\
 \hline
 \text{номер}&5&2&1&4&3&6
 \end{array}
\]
Исход --- упорядоченная четвёрка различных номеров из $D=\{1,\ldots,6\}$:
\[
 \Omega=\{(x_1,x_2,x_3,x_4):x_i\in D,\ x_i\ne x_j\text{ при }i\ne j\}.
\]
Можно выбирать последовательно:
$x_1\in D$, $x_2\in D\setminus\{x_1\}$,
$x_3\in D\setminus\{x_1,x_2\}$,
$x_4\in D\setminus\{x_1,x_2,x_3\}$.

Сочетания --- выбор нескольких элементов без учёта порядка:
из $\{\text{А},\text{Б},\text{В}\}$ по два получаем АБ, АВ, БВ; БА и АБ совпадают.
Для размещений порядок важен.
\[
 C_n^k=\frac{n!}{k!(n-k)!},\qquad A_n^k=\frac{n!}{(n-k)!},
\]
где $n$ --- число всех элементов, $k$ --- число выбираемых.
\[
 N=A_6^4=\frac{6!}{(6-4)!}=360,
 \qquad A=\{(1,2,3,4)\},\qquad N_A=1,
\]
\[
 P(A)=\frac1{A_6^4}=\frac1{360}.
\]
\editorial{Полное словесное условие на этой странице не записано.
Событие сохранено в виде четвёрки номеров; по указанной нумерации это буквы «ТИРЕ».
В одной строке исходника вместо $A_6^4$ ошибочно написано $C_6^4$; обозначение исправлено.}

\sourcepage{IMG_7623.jpeg}
\section*{Задача 2.23. Нумерованные шары}
В урне $10$ пронумерованных шаров. Их вынимают один за другим, пока
не будут вынуты все. Найти вероятность того, что номера идут по порядку.
\[
 \Omega=\{(a_1,\ldots,a_{10}):a_i\in\{1,\ldots,10\},\
 a_i\ne a_j\text{ при }i\ne j\}.
\]
Порядок важен, поэтому
\[
 N=A_{10}^{10}=\frac{10!}{0!}=10!.
\]
\[
 A=\{(1,2,3,4,5,6,7,8,9,10)\},\qquad N_A=1,
 \qquad P(A)=\frac1{10!}.
\]

\section*{Задача 2.29. Студенты и вагоны}
Десять студентов договорились о поездке за город, но не договорились о вагоне.
Каждый может сесть в любой из десяти вагонов. Какова вероятность,
что они окажутся в разных вагонах?

\sourcepage{IMG_7624.jpeg}
Исход $(x_1,\ldots,x_{10})$, где $x_i\in\{1,\ldots,10\}$ ---
номер вагона $i$-го студента. Порядок важен, повторения возможны:
\[
 N=10^{10}.
\]
Событию $A$ соответствуют исходы с попарно различными номерами вагонов:
\[
 N_A=A_{10}^{10}=10!,\qquad
 P(A)=\frac{N_A}{N}=\frac{10!}{10^{10}}.
\]

\section*{Задача 2.35. Карты}
Какова вероятность того, что среди четырёх карт, выбранных наудачу
из колоды в $52$ карты, ровно две окажутся трефовой масти?
\editorial{Записанное ниже решение целиком зачёркнуто. Оно не восстанавливается и не дополняется.}

\end{document}
